\begin{PSbox}[unbreakable]{obj}{Learning Objectives}

After completing this module, students will be able to:

\begin{enumerate}
\def\labelenumi{\arabic{enumi}.}
\tightlist
\item
  Define and compute joint PMF, PDF, and CDF
\item
  Derive marginal distributions from joint distributions
\item
  Compute conditional distributions
\item
  Test independence of two random variables
\item
  Compute conditional expectations
\item
  Visualize joint distributions in MATLAB
\end{enumerate}

\end{PSbox}

\PSsection{6.1}{Introduction: Why Study Joint Distributions?}

Real engineering systems involve \textbf{multiple interacting uncertain quantities}:

\begin{itemize}
\tightlist
\item
  Voltage AND current in a noisy circuit
\item
  Signal AND noise at a receiver
\item
  Temperature AND resistance of a sensor
\item
  I-channel AND Q-channel of a communication signal
\end{itemize}

Understanding the \textbf{joint behavior} tells us things marginals cannot — like correlation, dependence, and conditional behavior.

\PSsection{6.2}{Joint PMF (Discrete Case)}

\begin{PSbox}[unbreakable]{def}{Definition}

For discrete RVs $X,\, Y$:\PSdisplay{p_{X,Y}(x,y) = P(X = x, Y = y)}

\end{PSbox}

\PSsubsection{Properties}

\begin{enumerate}
\def\labelenumi{\arabic{enumi}.}
\tightlist
\item
  $p_{X,Y}(x,\,y) \ge 0$
\item
  $\sum_{x} \Sigma_{y} p_{X,Y}(x,\,y) = 1$
\end{enumerate}

\PSsubsection{Visualization: Probability Table}

\textbf{Example:} $X$ = number of retransmissions $(0,\,1,\,2),\, Y$ = acknowledgment delay category $(1 = \text{fast},\, 2 = \text{slow})$

{\def\LTcaptype{none} % do not increment counter
\begin{longtable}[c]{@{}cccc@{}}
\toprule\noalign{}
\rowcolor{PSnavy}& \PSth{$Y = 1$} & \PSth{$Y = 2$} & \PSth{$p_{X}(x)$} \\
\midrule\noalign{}
\endhead
\bottomrule\noalign{}
\endlastfoot
$X = 0$ & 0.30 & 0.10 & 0.40 \\
$X = 1$ & 0.15 & 0.20 & 0.35 \\
$X = 2$ & 0.05 & 0.20 & 0.25 \\
$p_{Y}(y)$ & 0.50 & 0.50 & 1.00 \\
\end{longtable}
}

\PSsection{6.3}{Joint PDF (Continuous Case)}

\begin{PSbox}[unbreakable]{def}{Definition}

For continuous RVs $X,\, Y$, the joint PDF $f_{X,Y}(x,\,y)$ satisfies:\PSdisplay{P((X,Y) \in A) = \iint_A f_{X,Y}(x,y) \, dx \, dy}

\end{PSbox}

\PSsubsection{Properties}

\begin{enumerate}
\def\labelenumi{\arabic{enumi}.}
\tightlist
\item
  $f_{X,Y}(x,\,y) \ge 0$
\item
  $\int \int f_{X,Y}(x,\,y)\,dx\,dy = 1$
\end{enumerate}

\PSsubsection{Visualization}

The joint PDF is a \textbf{surface} over the $(x,\,y)$ plane. Probability = volume under surface over a region.

\begin{PSbox}[unbreakable]{ex}{Engineering Example: Bivariate Gaussian}

Two correlated noise voltages $(X,\, Y)$ with correlation $\rho$:

\PSdisplay{f_{X,Y}(x,y) = \frac{1}{2\pi\sigma_X\sigma_Y\sqrt{1-\rho^2}} \exp\left(-\frac{1}{2(1-\rho^2)}\left[\frac{x^2}{\sigma_X^2} - \frac{2\rho xy}{\sigma_X\sigma_Y} + \frac{y^2}{\sigma_Y^2}\right]\right)}

\end{PSbox}

\PSsection{6.4}{Joint CDF}

\begin{PSbox}[unbreakable]{def}{Definition}

\PSdisplay{F_{X,Y}(x,y) = P(X \leq x, Y \leq y)}

\end{PSbox}

\PSsubsection{Properties}

\begin{enumerate}
\def\labelenumi{\arabic{enumi}.}
\tightlist
\item
  $F_{X,Y}(-\infty,\, y) = 0,\, F_{X,Y}(x,\, -\infty) = 0$
\item
  $F_{X,Y}(\infty,\, \infty) = 1$
\item
  Non-decreasing in both arguments
\item
  $\displaystyle f_{X,Y}(x,\,y) = \frac{\partial^{2}F_{X,Y}}{\partial x \partial y}$
\end{enumerate}

\PSsection{6.5}{Marginal Distributions}

\PSsubsection{Obtaining Marginals from Joint}

\textbf{Discrete:}\PSdisplay{p_X(x) = \sum_y p_{X,Y}(x,y), \quad p_Y(y) = \sum_x p_{X,Y}(x,y)}

\textbf{Continuous:}\PSdisplay{f_X(x) = \int_{-\infty}^{\infty} f_{X,Y}(x,y) \, dy, \quad f_Y(y) = \int_{-\infty}^{\infty} f_{X,Y}(x,y) \, dx}

\PSsubsection{Visual Interpretation}

The marginal PDF $f_{X}(x)$ is the \textbf{projection} (integral) of the joint surface onto the x-axis. Similarly $f_{Y}(y)$ projects onto the y-axis.

\begin{PSbox}[unbreakable]{ex}{Engineering Example}

Joint PDF: $f_{X,Y}(x,\,y) = 2e^{-x}e^{-2y},\, x \ge 0,\, y \ge 0$

Marginals:

\begin{itemize}
\tightlist
\item
  $\displaystyle f_{X}(x) = \int_{0}^{\infty} 2e^{-x}e^{-2y}\,dy = 2e^{-x} \cdot \left[\frac{1}{2}\right] = e^{-x} \to X \sim \mathrm{Exp}(1)$
\item
  $f_{Y}(y) = \int_{0}^{\infty} 2e^{-x}e^{-2y}\,dx = 2e^{-2y} \cdot [1] = 2e^{-2y} \to Y \sim \mathrm{Exp}(2)$
\end{itemize}

\end{PSbox}

\PSsection{6.6}{Conditional Distributions}

\PSsubsection{Discrete}

\PSdisplay{p_{X|Y}(x|y) = \frac{p_{X,Y}(x,y)}{p_Y(y)}}

\PSsubsection{Continuous}

\PSdisplay{f_{X|Y}(x|y) = \frac{f_{X,Y}(x,y)}{f_Y(y)}}

\PSsubsection{Interpretation}

The conditional distribution describes $X$ when we \textbf{know} $Y = y$. It’s a ``slice'' of the joint distribution at a fixed $y$ value, renormalized to integrate to 1.

\begin{PSbox}[unbreakable]{ex}{Engineering Example: Signal Given Channel State}

Received power $Y$ given channel state $X$:

\begin{itemize}
\tightlist
\item
  If channel is good $(X = 1)$: $Y \sim N(10,\, 1)$
\item
  If channel is poor $(X = 0)$: $Y \sim N(2,\, 4)$
\item
  $P(X = 1) = 0.7,\, P(X = 0) = 0.3$
\end{itemize}

Conditional PDFs describe receiver behavior in each state.

\end{PSbox}

\PSsection{6.7}{Independence of Two Random Variables}

\begin{PSbox}[unbreakable]{def}{Definition}

$X$ and $Y$ are independent if and only if:

\PSdisplay{f_{X,Y}(x,y) = f_X(x) \cdot f_Y(y) \quad \text{for all } x, y}

(Or equivalently for PMFs: $p_{X,Y}(x,\,y) = p_{X}(x) \cdot p_{Y}(y)$)

\end{PSbox}

\PSsubsection{Consequences of Independence}

\begin{itemize}
\tightlist
\item
  $f_{X \mid Y}(x \mid y) = f_{X}(x)$ (knowing $Y$ doesn’t change $X$)
\item
  $E[XY] = E[X]E[Y]$
\item
  $\operatorname{Var}(X + Y) = \operatorname{Var}(X) + \operatorname{Var}(Y)$
\end{itemize}

\PSsubsection{Testing Independence}

Check: Does the joint factor into a product of a function of $x$ alone and $y$ alone?

\begin{itemize}
\tightlist
\item
  $f_{X,Y}(x,\,y) = 2e^{-x} \cdot e^{-2y} = [e^{-x}] \cdot [2e^{-2y}]$ \PSyes{} Independent!
\item
  $f_{X,Y}(x,\,y) = x + y$ for $0 \le x,\,y \le 1$ \ensuremath{\to} cannot factor \ensuremath{\to} NOT independent
\end{itemize}

\begin{PSbox}[unbreakable]{ex}{Engineering Example: Signal and Independent Noise}

Transmitted signal $S$ and channel noise $N$ are independent by physical assumption:\PSbr{}$f_{S,N}(s,\,n) = f_{S}(s) \cdot f_{N}(n)$

This independence assumption is fundamental to communication system design.

\end{PSbox}

\PSsection{6.8}{Conditional Expectation}

\begin{PSbox}[unbreakable]{def}{Definition}

\PSdisplay{E[X|Y=y] = \begin{cases} \sum_x x \cdot p_{X|Y}(x|y) & \text{discrete} \\ \int_{-\infty}^{\infty} x \cdot f_{X|Y}(x|y) \, dx & \text{continuous} \end{cases}}

\end{PSbox}

\PSsubsection{Interpretation}

$E[X \mid Y = y]$ is the \textbf{best prediction} of $X$ given that we observed $Y = y$ (in the mean-square sense).

\begin{PSbox}[unbreakable]{ex}{Engineering Example: Temperature-Dependent Resistance}

A resistor’s value $R$ depends on temperature $T$. If $R \mid T \sim N(1000 + 0.5T,\, 4)$:

$E[R \mid T = 50^{\circ}C] = 1000 + 0.5(50) = 1025\,\Omega$

Given temperature, we can predict resistance.

\end{PSbox}

\PSsection{6.9}{MATLAB Examples}

\begin{PSbox}[unbreakable]{ex}{Example 1: Joint PMF Visualization}

\tcbset{PScodemode/.style={unbreakable}}

\begin{Shaded}
\begin{Highlighting}[]
\CommentTok{\%\% Joint PMF: Retransmissions vs Delay Category}
\VariableTok{pXY} \OperatorTok{=}\NormalTok{ [}\FloatTok{0.30} \FloatTok{0.10}\OperatorTok{;} \FloatTok{0.15} \FloatTok{0.20}\OperatorTok{;} \FloatTok{0.05} \FloatTok{0.20}\NormalTok{]}\OperatorTok{;}
\VariableTok{x} \OperatorTok{=} \FloatTok{0}\OperatorTok{:}\FloatTok{2}\OperatorTok{;} \VariableTok{y} \OperatorTok{=} \FloatTok{1}\OperatorTok{:}\FloatTok{2}\OperatorTok{;}

\VariableTok{figure}\OperatorTok{;}
\VariableTok{bar3}\NormalTok{(}\VariableTok{pXY}\NormalTok{)}\OperatorTok{;}
\VariableTok{xlabel}\NormalTok{(}\SpecialStringTok{\textquotesingle{}Delay Category Y\textquotesingle{}}\NormalTok{)}\OperatorTok{;} \VariableTok{ylabel}\NormalTok{(}\SpecialStringTok{\textquotesingle{}Retransmissions X\textquotesingle{}}\NormalTok{)}\OperatorTok{;}
\VariableTok{zlabel}\NormalTok{(}\SpecialStringTok{\textquotesingle{}P(X=x, Y=y)\textquotesingle{}}\NormalTok{)}\OperatorTok{;}
\VariableTok{title}\NormalTok{(}\SpecialStringTok{\textquotesingle{}Joint PMF\textquotesingle{}}\NormalTok{)}\OperatorTok{;}

\CommentTok{\% Marginals}
\VariableTok{pX} \OperatorTok{=} \VariableTok{sum}\NormalTok{(}\VariableTok{pXY}\OperatorTok{,} \FloatTok{2}\NormalTok{)}\OperatorTok{\textquotesingle{};}   \CommentTok{\% Sum over Y}
\VariableTok{pY} \OperatorTok{=} \VariableTok{sum}\NormalTok{(}\VariableTok{pXY}\OperatorTok{,} \FloatTok{1}\NormalTok{)}\OperatorTok{;}    \CommentTok{\% Sum over X}
\VariableTok{fprintf}\NormalTok{(}\SpecialStringTok{\textquotesingle{}Marginal of X: [\%.2f \%.2f \%.2f]\textbackslash{}n\textquotesingle{}}\OperatorTok{,} \VariableTok{pX}\NormalTok{)}\OperatorTok{;}
\VariableTok{fprintf}\NormalTok{(}\SpecialStringTok{\textquotesingle{}Marginal of Y: [\%.2f \%.2f]\textbackslash{}n\textquotesingle{}}\OperatorTok{,} \VariableTok{pY}\NormalTok{)}\OperatorTok{;}
\end{Highlighting}
\end{Shaded}

\end{PSbox}

\begin{PSbox}[unbreakable]{ex}{Example 2: Bivariate Gaussian Visualization}

\tcbset{PScodemode/.style={unbreakable}}

\begin{Shaded}
\begin{Highlighting}[]
\CommentTok{\%\% Joint Gaussian PDF with Correlation}
\VariableTok{mu} \OperatorTok{=}\NormalTok{ [}\FloatTok{0} \FloatTok{0}\NormalTok{]}\OperatorTok{;} \VariableTok{sigma\_x} \OperatorTok{=} \FloatTok{1}\OperatorTok{;} \VariableTok{sigma\_y} \OperatorTok{=} \FloatTok{1.5}\OperatorTok{;} \VariableTok{rho} \OperatorTok{=} \FloatTok{0.7}\OperatorTok{;}
\VariableTok{Sigma} \OperatorTok{=}\NormalTok{ [}\VariableTok{sigma\_x}\OperatorTok{\^{}}\FloatTok{2}\OperatorTok{,} \VariableTok{rho}\OperatorTok{*}\VariableTok{sigma\_x}\OperatorTok{*}\VariableTok{sigma\_y}\OperatorTok{;} \VariableTok{rho}\OperatorTok{*}\VariableTok{sigma\_x}\OperatorTok{*}\VariableTok{sigma\_y}\OperatorTok{,} \VariableTok{sigma\_y}\OperatorTok{\^{}}\FloatTok{2}\NormalTok{]}\OperatorTok{;}

\NormalTok{[}\VariableTok{X}\OperatorTok{,} \VariableTok{Y}\NormalTok{] }\OperatorTok{=} \VariableTok{meshgrid}\NormalTok{(}\VariableTok{linspace}\NormalTok{(}\OperatorTok{{-}}\FloatTok{4}\OperatorTok{,} \FloatTok{4}\OperatorTok{,} \FloatTok{100}\NormalTok{)}\OperatorTok{,} \VariableTok{linspace}\NormalTok{(}\OperatorTok{{-}}\FloatTok{5}\OperatorTok{,} \FloatTok{5}\OperatorTok{,} \FloatTok{100}\NormalTok{))}\OperatorTok{;}
\VariableTok{XY} \OperatorTok{=}\NormalTok{ [}\VariableTok{X}\NormalTok{(}\OperatorTok{:}\NormalTok{) }\VariableTok{Y}\NormalTok{(}\OperatorTok{:}\NormalTok{)]}\OperatorTok{;}
\VariableTok{Z} \OperatorTok{=} \VariableTok{mvnpdf}\NormalTok{(}\VariableTok{XY}\OperatorTok{,} \VariableTok{mu}\OperatorTok{,} \VariableTok{Sigma}\NormalTok{)}\OperatorTok{;}
\VariableTok{Z} \OperatorTok{=} \VariableTok{reshape}\NormalTok{(}\VariableTok{Z}\OperatorTok{,} \VariableTok{size}\NormalTok{(}\VariableTok{X}\NormalTok{))}\OperatorTok{;}

\VariableTok{figure}\OperatorTok{;}
\VariableTok{subplot}\NormalTok{(}\FloatTok{1}\OperatorTok{,}\FloatTok{2}\OperatorTok{,}\FloatTok{1}\NormalTok{)}\OperatorTok{;}
\VariableTok{surf}\NormalTok{(}\VariableTok{X}\OperatorTok{,} \VariableTok{Y}\OperatorTok{,} \VariableTok{Z}\OperatorTok{,} \SpecialStringTok{\textquotesingle{}EdgeColor\textquotesingle{}}\OperatorTok{,} \SpecialStringTok{\textquotesingle{}none\textquotesingle{}}\NormalTok{)}\OperatorTok{;} \VariableTok{view}\NormalTok{(}\FloatTok{30}\OperatorTok{,} \FloatTok{40}\NormalTok{)}\OperatorTok{;}
\VariableTok{xlabel}\NormalTok{(}\SpecialStringTok{\textquotesingle{}X\textquotesingle{}}\NormalTok{)}\OperatorTok{;} \VariableTok{ylabel}\NormalTok{(}\SpecialStringTok{\textquotesingle{}Y\textquotesingle{}}\NormalTok{)}\OperatorTok{;} \VariableTok{zlabel}\NormalTok{(}\SpecialStringTok{\textquotesingle{}f\_\{X,Y\}(x,y)\textquotesingle{}}\NormalTok{)}\OperatorTok{;}
\VariableTok{title}\NormalTok{(}\VariableTok{sprintf}\NormalTok{(}\SpecialStringTok{\textquotesingle{}Joint Gaussian (ρ = \%.1f)\textquotesingle{}}\OperatorTok{,} \VariableTok{rho}\NormalTok{))}\OperatorTok{;}

\VariableTok{subplot}\NormalTok{(}\FloatTok{1}\OperatorTok{,}\FloatTok{2}\OperatorTok{,}\FloatTok{2}\NormalTok{)}\OperatorTok{;}
\VariableTok{contour}\NormalTok{(}\VariableTok{X}\OperatorTok{,} \VariableTok{Y}\OperatorTok{,} \VariableTok{Z}\OperatorTok{,} \FloatTok{20}\NormalTok{)}\OperatorTok{;}
\VariableTok{xlabel}\NormalTok{(}\SpecialStringTok{\textquotesingle{}X\textquotesingle{}}\NormalTok{)}\OperatorTok{;} \VariableTok{ylabel}\NormalTok{(}\SpecialStringTok{\textquotesingle{}Y\textquotesingle{}}\NormalTok{)}\OperatorTok{;}
\VariableTok{title}\NormalTok{(}\SpecialStringTok{\textquotesingle{}Contour Plot\textquotesingle{}}\NormalTok{)}\OperatorTok{;}
\VariableTok{axis} \VariableTok{equal}\OperatorTok{;}
\end{Highlighting}
\end{Shaded}

\end{PSbox}

\begin{PSbox}[unbreakable]{ex}{Example 3: Testing Independence}

\tcbset{PScodemode/.style={unbreakable}}

\begin{Shaded}
\begin{Highlighting}[]
\CommentTok{\%\% Independence Check via Simulation}
\VariableTok{N} \OperatorTok{=} \FloatTok{100000}\OperatorTok{;}

\CommentTok{\% Case 1: Independent (signal and noise)}
\VariableTok{S} \OperatorTok{=} \VariableTok{randn}\NormalTok{(}\FloatTok{1}\OperatorTok{,} \VariableTok{N}\NormalTok{)}\OperatorTok{;}               \CommentTok{\% Signal}
\VariableTok{Noise} \OperatorTok{=} \FloatTok{0.5}\OperatorTok{*}\VariableTok{randn}\NormalTok{(}\FloatTok{1}\OperatorTok{,} \VariableTok{N}\NormalTok{)}\OperatorTok{;}       \CommentTok{\% Independent noise}
\VariableTok{R} \OperatorTok{=} \VariableTok{S} \OperatorTok{+} \VariableTok{Noise}\OperatorTok{;}                  \CommentTok{\% Received}

\CommentTok{\% Case 2: Dependent (voltage and current through same resistor)}
\VariableTok{V} \OperatorTok{=} \VariableTok{randn}\NormalTok{(}\FloatTok{1}\OperatorTok{,} \VariableTok{N}\NormalTok{)}\OperatorTok{;}               \CommentTok{\% Voltage}
\VariableTok{I} \OperatorTok{=} \VariableTok{V} \OperatorTok{/} \FloatTok{50} \OperatorTok{+} \FloatTok{0.01}\OperatorTok{*}\VariableTok{randn}\NormalTok{(}\FloatTok{1}\OperatorTok{,}\VariableTok{N}\NormalTok{)}\OperatorTok{;} \CommentTok{\% Current (dependent on V!)}

\VariableTok{fprintf}\NormalTok{(}\SpecialStringTok{\textquotesingle{}Independent (S, Noise): corr = \%.4f\textbackslash{}n\textquotesingle{}}\OperatorTok{,} \VariableTok{corr}\NormalTok{(}\VariableTok{S}\OperatorTok{\textquotesingle{},} \VariableTok{Noise}\OperatorTok{\textquotesingle{}}\NormalTok{))}\OperatorTok{;}
\VariableTok{fprintf}\NormalTok{(}\SpecialStringTok{\textquotesingle{}Dependent (V, I): corr = \%.4f\textbackslash{}n\textquotesingle{}}\OperatorTok{,} \VariableTok{corr}\NormalTok{(}\VariableTok{V}\OperatorTok{\textquotesingle{},} \VariableTok{I}\OperatorTok{\textquotesingle{}}\NormalTok{))}\OperatorTok{;}
\end{Highlighting}
\end{Shaded}

\end{PSbox}

\begin{PSbox}[unbreakable]{ex}{Example 4: Conditional Distribution}

\tcbset{PScodemode/.style={unbreakable}}

\begin{Shaded}
\begin{Highlighting}[]
\CommentTok{\%\% Conditional PDF: Slice of Bivariate Gaussian}
\VariableTok{rho} \OperatorTok{=} \FloatTok{0.8}\OperatorTok{;} \VariableTok{N} \OperatorTok{=} \FloatTok{100000}\OperatorTok{;}
\VariableTok{X} \OperatorTok{=} \VariableTok{randn}\NormalTok{(}\FloatTok{1}\OperatorTok{,} \VariableTok{N}\NormalTok{)}\OperatorTok{;}
\VariableTok{Y} \OperatorTok{=} \VariableTok{rho}\OperatorTok{*}\VariableTok{X} \OperatorTok{+} \VariableTok{sqrt}\NormalTok{(}\FloatTok{1}\OperatorTok{{-}}\VariableTok{rho}\OperatorTok{\^{}}\FloatTok{2}\NormalTok{)}\OperatorTok{*}\VariableTok{randn}\NormalTok{(}\FloatTok{1}\OperatorTok{,} \VariableTok{N}\NormalTok{)}\OperatorTok{;}  \CommentTok{\% Y|X \textasciitilde{} N(rho*x, 1{-}rho\^{}2)}

\CommentTok{\% Conditional: Y given X ≈ 1 (take samples where 0.9 \textless{} X \textless{} 1.1)}
\VariableTok{mask} \OperatorTok{=}\NormalTok{ (}\VariableTok{X} \OperatorTok{\textgreater{}} \FloatTok{0.9}\NormalTok{) }\OperatorTok{\&}\NormalTok{ (}\VariableTok{X} \OperatorTok{\textless{}} \FloatTok{1.1}\NormalTok{)}\OperatorTok{;}
\VariableTok{Y\_cond} \OperatorTok{=} \VariableTok{Y}\NormalTok{(}\VariableTok{mask}\NormalTok{)}\OperatorTok{;}

\VariableTok{figure}\OperatorTok{;}
\VariableTok{histogram}\NormalTok{(}\VariableTok{Y\_cond}\OperatorTok{,} \FloatTok{50}\OperatorTok{,} \SpecialStringTok{\textquotesingle{}Normalization\textquotesingle{}}\OperatorTok{,} \SpecialStringTok{\textquotesingle{}pdf\textquotesingle{}}\NormalTok{)}\OperatorTok{;}
\VariableTok{hold} \VariableTok{on}\OperatorTok{;}
\VariableTok{y} \OperatorTok{=} \VariableTok{linspace}\NormalTok{(}\OperatorTok{{-}}\FloatTok{3}\OperatorTok{,} \FloatTok{4}\OperatorTok{,} \FloatTok{200}\NormalTok{)}\OperatorTok{;}
\VariableTok{plot}\NormalTok{(}\VariableTok{y}\OperatorTok{,} \VariableTok{normpdf}\NormalTok{(}\VariableTok{y}\OperatorTok{,} \VariableTok{rho}\OperatorTok{*}\FloatTok{1}\OperatorTok{,} \VariableTok{sqrt}\NormalTok{(}\FloatTok{1}\OperatorTok{{-}}\VariableTok{rho}\OperatorTok{\^{}}\FloatTok{2}\NormalTok{))}\OperatorTok{,} \SpecialStringTok{\textquotesingle{}r\textquotesingle{}}\OperatorTok{,} \SpecialStringTok{\textquotesingle{}LineWidth\textquotesingle{}}\OperatorTok{,} \FloatTok{2}\NormalTok{)}\OperatorTok{;}
\VariableTok{xlabel}\NormalTok{(}\SpecialStringTok{\textquotesingle{}Y\textquotesingle{}}\NormalTok{)}\OperatorTok{;} \VariableTok{ylabel}\NormalTok{(}\SpecialStringTok{\textquotesingle{}PDF\textquotesingle{}}\NormalTok{)}\OperatorTok{;}
\VariableTok{title}\NormalTok{(}\SpecialStringTok{\textquotesingle{}f\_\{Y|X\}(y|X≈1) for Bivariate Gaussian\textquotesingle{}}\NormalTok{)}\OperatorTok{;}
\VariableTok{legend}\NormalTok{(}\SpecialStringTok{\textquotesingle{}Empirical\textquotesingle{}}\OperatorTok{,} \VariableTok{sprintf}\NormalTok{(}\SpecialStringTok{\textquotesingle{}N(\%.1f, \%.2f)\textquotesingle{}}\OperatorTok{,} \VariableTok{rho}\OperatorTok{,} \FloatTok{1}\OperatorTok{{-}}\VariableTok{rho}\OperatorTok{\^{}}\FloatTok{2}\NormalTok{))}\OperatorTok{;}
\end{Highlighting}
\end{Shaded}

\end{PSbox}

\PSsection{6.10}{Practice Problems}

\begin{PSbox}[unbreakable]{prob}{Problem 1}

Joint PDF: $f_{X,Y}(x,\,y) = 6(1-y)$ for $0 \le x \le y \le 1$, zero elsewhere.

\begin{enumerate}
\def\labelenumi{(\alph{enumi})}
\tightlist
\item
  Verify normalization. (b) Find marginals. (c) Are X,Y independent? (d) Find $P(X < 0.5,\, Y < 0.5)$.
\end{enumerate}

\PSsolution{} 

\begin{enumerate}
\def\labelenumi{(\alph{enumi})}
\tightlist
\item
  $\int_{0}^{1} \int_{0}^{y} 6(1-y)\,dx\,dy = \int_{0}^{1} 6y(1-y)\,dy = 6\left[\frac{y^{2}}{2} - \frac{y^{3}}{3}\right]_{0}^{1} = 6\left(\frac{1}{2} - \frac{1}{3}\right) = 1$ \PSyes{}
\item
  $f_{X}(x) = \int_{x}^{1} 6(1-y)\,dy = 6\left[\left(1-\frac{y^{2}}{2}\right) - (y)\right]_{x}^{1} = 3(1-x)^{2}$; $f_{Y}(y) = 6y(1-y)$
\item
  $f_{X,Y} \ne f_{X} \cdot f_{Y}$ \ensuremath{\to} NOT independent (also, the support is triangular, not rectangular)
\item
  $\displaystyle P(X < 0.5,\, Y < 0.5) = \int_{0}^{0.5} \int_{0}^{y} 6(1-y)\,dx\,dy = \int_{0}^{0.5} 6y(1-y)\,dy = 6\left[\frac{y^{2}}{2} - \frac{y^{3}}{3}\right]_{0}^{0.5} = 0.500$
\end{enumerate}

\end{PSbox}

\begin{PSbox}[unbreakable]{prob}{Problem 2}

Two sensors measure the same signal with independent noise. $X = S + N_{1},\, Y = S + N_{2}$ where $S = 5\,\mathrm{V},\, N_{1} \sim N(0,\,1),\, N_{2} \sim N(0,\,4)$. Are $X$ and $Y$ independent? Find $E[X],\, E[Y]$, and explain why $\operatorname{Cov}(X,\,Y) \ne 0$.

\PSsolution{} $X$ and $Y$ are NOT independent because they share $S$ (though noise is independent). $E[X] = E[Y] = 5$. $\operatorname{Cov}(X,\,Y) = \operatorname{Cov}(S+N_{1},\, S+N_{2}) = \operatorname{Var}(S) + \operatorname{Cov}(S,\,N_{1}) + \operatorname{Cov}(N_{2},\,S) + \operatorname{Cov}(N_{1},\,N_{2})$. Since $S$ is constant, $\operatorname{Var}(S) = 0$, so if $S$ were random they’d be correlated. With $S$ fixed, $X$ and $Y$ are actually independent in this case. If $S$ is random with $\operatorname{Var}(S) = \sigma_{S}^{2}$, then $\operatorname{Cov}(X,\,Y) = \sigma_{S}^{2}$.

\end{PSbox}

\begin{PSbox}[unbreakable]{prob}{Problem 3}

Given joint PMF in the table of Section 6.2, find (a) $P(X \le 1,\, Y = 2)$, (b) $f_{X \mid Y}(x \mid Y = 1)$, (c) $E[X \mid Y = 1]$.

\PSsolution{} 

\begin{enumerate}
\def\labelenumi{(\alph{enumi})}
\tightlist
\item
  $P(X \le 1,\, Y = 2) = p(0,\,2) + p(1,\,2) = 0.10 + 0.20 = 0.30$
\item
  $p_{X \mid Y}(x \mid 1) = \frac{p(x,\,1)}{p_{Y}(1)}$: $p(0 \mid 1) = \frac{0.30}{0.50} = 0.60,\, p(1 \mid 1) = \frac{0.15}{0.50} = 0.30,\, p(2 \mid 1) = \frac{0.05}{0.50} = 0.10$
\item
  $E[X \mid Y = 1] = 0(0.60) + 1(0.30) + 2(0.10) = 0.50$
\end{enumerate}

\emph{Next Module: {Module 7 — Functions of Two Random Variables}}

\end{PSbox}
